\documentclass[10pt,a4paper]{amsart}
\usepackage[margin=19mm]{geometry}
\usepackage{amsmath,amssymb,mathtools}
\usepackage[scr=rsfs,cal=euler]{mathalfa}
\usepackage{xfrac}
\usepackage[unicode,hidelinks,bookmarks=false]{hyperref}
\newcommand{\ie}{i.e.\ }
\newcommand{\cf}{cf.\ }
\newcommand{\resp}{respectively\ }
\newcommand{\Subsection}{\subsection}
\providecommand{\nobreakdash}{\nobreak\hbox{-}}
\setlength{\emergencystretch}{2em}
\multlinegap0pt
\usepackage{amssymb}
\newtheorem{theorem}{Theorem}
\newtheorem{lemma}{Lemma}
\title{\textbf{A generalization of the Erdős-Sierpiński conjecture}}
\author{Amirali Fatehizadeh}
\date{Source: arXiv:2605.21524v1 (CC BY 4.0)}
\begin{document}
\maketitle

\noindent\emph{Source and licence.} \textbf{A generalization of the Erdős-Sierpiński conjecture}. Version arXiv:2605.21524v1 (CC BY 4.0), distributed by arXiv under the Creative Commons Attribution 4.0 International licence, http://creativecommons.org/licenses/by/4.0/, as recorded in the arXiv licence metadata for this exact version and shown on the abstract page https://arxiv.org/abs/2605.21524v1. Full licence notice: \texttt{licenses/arxiv-2605.21524-CC-BY-4.0.txt}.\par\smallskip

\noindent\textit{Editorial scope.} Counted unit: the article's theorem theorem (source0.tex lines 177--180). The article's complete original proof of it is delivered with it (source0.tex lines 182--286, one \texttt{proof} environment of 105 lines). No mathematics is written, repaired or re-derived: the statement and the proof are the article's own lines, in the article's own notation, and no retained line is substituted or re-flowed.  No further source-local context is needed: every cross-reference of every retained line resolves inside the two delivered spans.  Nothing was omitted from any retained span, so the package carries no omission record.  No retained line contains a citation, so the wrapper carries no bibliography and no bibliography entry.  No retained line contains a figure environment, an image-inclusion command or drawing code, so no image is delivered and the package declares no asset dependency.  Editorial formatting only, in addition: an AMS \texttt{amsart} preamble that restates the article's own declarations and macros used by the retained lines, namely the environments \texttt{theorem}, \texttt{lemma} on separate counters, together with the standard packages those lines need.  The retained environments print in this document as Theorem 1, Lemma 1 in order of appearance, whereas the article prints the same items with the numbering of its own full document.  The complete native source of the article as distributed in the arXiv e-print for this exact version is bundled unchanged as \texttt{sources/201068/source0.tex}.
\begin{theorem}
Let $k > 1$ be an integer number. For sufficiently large $x$, we have
\[ A_k(x) \ll_k \frac{x}{\sqrt{\log \log \log x}}. \]
\end{theorem}

\begin{proof}
Let $n \in A_k$. By dividing both sides of the main equation by $n(n + 1)$ and applying the logarithmic function, the equation is algebraically equivalent to the following form:
\[ g(n + 1) - g(n) = \log k - \log \left(1 + \frac{1}{n}\right). \]

To control the behavior of the function $g$, we truncate it using two constant parameters $y \geq 2$ and $r \geq 1$. We define the truncated function as $h_{y,r}(n) := \sum_{p \leq y} \ell\left(p^{\min{(v_p(n), r)}}\right)$ and its consecutive difference as $D_{y,r}(n) := h_{y,r}(n + 1) - h_{y,r}(n)$.

For a real number $x > 0$ and an error parameter $\epsilon > 0$, we construct the following sets in the interval $[1, x]$:
\begin{enumerate}
    \item $E_0(x, \epsilon) := \left\{n \leq x : \log \left(1 + \frac{1}{n}\right) > \epsilon \right\}$
    \item $E_1(x, y, \epsilon) := \{n \leq x : g_{>y}(n) > \epsilon \lor g_{>y}(n + 1) > \epsilon\}$
    \item $E_2(x, y, r) := \{n \leq x : \exists p \leq y, v_p(n) > r \lor v_p(n + 1) > r\}$
    \item $S(x, y, r, \epsilon) := \{n \leq x : |D_{y,r}(n) - \log k| \leq 3\epsilon\}$
\end{enumerate}

Suppose $n \leq x$ is a solution to the equation that does not belong to any of the sets $E_0, E_1$, and $E_2$. Since $n \notin E_2$, the decomposition $g(n) = h_{y,r}(n) + g_{>y}(n)$ is valid for both $n$ and $n + 1$. By substituting this into the main equation and applying the triangle inequality, the conditions $n \notin E_0$ and $n \notin E_1$ imply that $|D_{y,r}(n) - \log k| \leq 3\epsilon$. Therefore, the set inclusion $A_k \cap [1, x] \subseteq E_0 \cup E_1 \cup E_2 \cup S$ strictly holds, and we have:
\begin{equation}
A_k(x) \leq \text{\#}E_0(x, \epsilon) + \text{\#}E_1(x, y, \epsilon) + \text{\#}E_2(x, y, r) + \text{\#}S(x, y, r, \epsilon). \label{eq:1}
\end{equation}

Now, we calculate an explicit upper bound for each of these error sets.

For the bound of $E_0$, condition $\log \left(1 + 1/n\right) > \epsilon$ implies $n < 1/\epsilon$. Thus, $\text{\#}E_0(x, \epsilon) \leq \frac{1}{\epsilon}$.

For the bound of $E_1$, considering the non-negativity of the terms $f_y(n) := g_{>y}(n) = \sum_{p > y} \ell(p^{v_p(n)})$, for any prime $p$ and integer $a \geq 1$, we have:
\[ \ell(p^a) = \log \left( \frac{1 - p^{-(a+1)}}{1 - p^{-1}} \right) \leq \log \left(\frac{1}{1 - p^{-1}}\right) = \log \left(\frac{p}{p - 1}\right) \leq \frac{2}{p}. \]

On the other hand, $\text{\#}\{n \leq x + 1 : v_p(n) = a\} \leq \lfloor \frac{x+1}{p^a} \rfloor \leq \frac{x+1}{p^a}$. Consequently, we obtain:
\[ \sum_{n \leq x+1} f_y(n) \leq (x + 1) \sum_{p > y} \sum_{a \geq 1} \frac{2}{p^{a+1}} = 2(x + 1) \sum_{p > y} \frac{1}{p(p - 1)}. \]

Now, using the Prime Number Theorem and partial summation, we have $\sum_{n \leq x+1} f_y(n) \ll \frac{x}{y \log y}$. Applying Markov's inequality to this non-negative function and considering the union bound for the two events associated with $n$ and $n + 1$, we get:
\[ \text{\#}E_1(x, y, \epsilon) \ll \frac{x}{\epsilon \cdot y \log y}. \]

For the bound of $E_2$, the number of integers up to $x$ that are divisible by $p^{r+1}$ is at most $\lfloor \frac{x}{p^{r+1}} \rfloor$. Using the integral test for the series over all integers $n \geq 2$, we have:
\[ \sum_{p \leq y} \frac{1}{p^{r+1}} \leq \sum_{n=2}^{\infty} \frac{1}{n^{r+1}} \leq \frac{1}{2^{r+1}} + \int_{2}^{\infty} \frac{dt}{t^{r+1}} = \frac{1}{2^{r+1}} + \frac{1}{r \cdot 2^r} \ll 2^{-r}. \]
Consequently, $\text{\#}E_2(x, y, r) \ll x \cdot 2^{-r}$.

Now, the function $D_{y,r}(n)$ over the natural numbers is strictly periodic, with a period of $M_{y,r} := \prod_{p \leq y} p^r$. For sufficiently large $x$, the parameter $y$ will also be sufficiently large. Thus, by the Prime Number Theorem and Chebyshev's inequality, we can write $\log M_{y,r} = r\theta(y) < 1.02ry$, and consequently $M_{y,r} \leq \exp(1.02ry)$.

We divide the interval $[1, x]$ into $m = \lfloor x / M_{y,r} \rfloor$ full periods and a remaining part of length strictly less than $M_{y,r}$. We equip the finite space $\mathbb{Z}/M_{y,r}\mathbb{Z}$ with a uniform probability measure. By the Chinese Remainder Theorem, the following ring isomorphism holds:
\[ \mathbb{Z}/M_{y,r}\mathbb{Z} \cong \prod_{p \leq y} \mathbb{Z}/p^r\mathbb{Z}. \]

Under this isomorphism, the uniform measure on the left-hand finite space is transferred to the product of uniform measures on the right-hand spaces. Consequently, if the random variable $N$ follows a uniform distribution on $\mathbb{Z}/M_{y,r}\mathbb{Z}$, its components $N_p := N \pmod{p^r}$ will be mutually independent random variables uniformly distributed on $\mathbb{Z}/p^r\mathbb{Z}$. Therefore, the distribution of $D_{y,r}(n)$ over each full period is exactly identical to the distribution of the following random variable:
\[ W_{y,r} := \sum_{p \leq y} Z_{p,r}, \quad Z_{p,r} := \ell\left(p^{\min(v_p(N_p + 1), r)}\right) - \ell\left(p^{\min(v_p(N_p), r)}\right). \]

Since the remaining part introduces an error bounded by its own length, $M_{y,r}$, we obtain:
\begin{equation}
\text{\#}S(x, y, r, \epsilon) \leq x \cdot \mathbb{P}(|W_{y,r} - \log k| \leq 3\epsilon) + \exp(1.02ry). \label{eq:2}
\end{equation}

Now, to apply Petrov's theorem, we first formulate the structure of the variables $Z_{p,r}$ in the form of the following lemma.

\begin{lemma}
Suppose $p \geq 5$ and the length of a closed interval $L$ is $L$. If $L < \ell(p)$, then for every $r \geq 1$ we have:
\[ Q_L(Z_{p,r}) = 1 - \frac{2}{p}. \]
\end{lemma}

\begin{proof}
The variables $N_p$ and $N_p + 1$ are consecutive; thus, they cannot simultaneously be multiples of $p$. This restricts the support of the variable $Z_{p,r}$ to three disjoint clusters:
\begin{enumerate}
    \item \textbf{Zero cluster:} $\mathbb{P}(Z_{p,r} = 0) = \mathbb{P}(p \nmid N_p(N_p + 1)) = 1 - 2/p$.
    \item \textbf{Positive cluster:} If $p \mid N_p + 1$. The total probability mass of this cluster is $1/p$, and its values fall within the interval $[\ell(p), \ell(p^r)]$; we have $\mathbb{P}(Z_{p,r} > 0) = 1/p$.
    \item \textbf{Negative cluster:} If $p \mid N_p$. Similar to the positive cluster, the total mass is $1/p$, and it takes values within the interval $[-\ell(p^r), -\ell(p)]$; we have $\mathbb{P}(Z_{p,r} < 0) = 1/p$.
\end{enumerate}

Therefore, we have $\text{supp}(Z_{p,r}) \subset \{0\} \cup \{\pm \ell(p^a)\}_{a=1}^r$. Since the minimum geometric distance between zero and any non-zero point is $\ell(p)$, and $L < \ell(p)$, no closed interval of length $L$ can simultaneously cover more than one component of the support. Thus, the maximum mass occurs exactly at zero. This completes the proof of Lemma 3.2.
\end{proof}

We now define the set of effective primes as $\mathcal{P}^* := \{p \mid 5 \leq p \leq \min(y, \frac{1}{13\epsilon})\}$. For any $p$ in this set, utilizing the bound $\log(1 + x) \geq x/2$, we have:
\[ \ell(p) = \log \left(1 + \frac{1}{p}\right) \geq \frac{1}{2p} \geq 6.5\epsilon > 6\epsilon. \]

Given the conditions of the preceding lemma, by setting the scale $L = 6\epsilon$, for all $p \in \mathcal{P}^*$ we obtain $Q_{6\epsilon}(Z_{p,r}) = 1 - 2/p$.

We decompose the independent random variable $W_{y,r}$ as $W_{y,r} = S^* + S'^*$, where $S^* := \sum_{p \in \mathcal{P}^*} Z_{p,r}$. Since the convolution with an independent random variable can only decrease the concentration function, we have $Q_{6\epsilon}(W_{y,r}) \leq Q_{6\epsilon}(S^*)$. Now, we apply Petrov's inequality (discussed in the previous section) on $S^*$ with both global and local scales set to $L = L_p = 6\epsilon$. By substitution, we obtain:
\[ Q_{6\epsilon}(W_{y,r}) \leq A \cdot 6\epsilon \cdot \left( \sum_{p \in \mathcal{P}^*} (6\epsilon)^2 (2/p) \right)^{-1/2}. \]

Upon simplification, the following inequality is deduced:
\begin{equation}
\mathbb{P}(|W_{y,r} - \log k| \leq 3\epsilon) \leq A \left( \sum_{p \in \mathcal{P}^*} \left(\frac{2}{p}\right) \right)^{-1/2}. \label{eq:3}
\end{equation}

Combining inequalities (\ref{eq:1}), (\ref{eq:2}), and (\ref{eq:3}), for all sufficiently large $x$, we have:
\[ \frac{A_k(x)}{x} \leq \frac{1}{x\epsilon} + \frac{C_1}{\epsilon \cdot y \log y} + C_2(2^{-r}) + \frac{\exp(1.02ry)}{x} + A \left( \sum_{p \in \mathcal{P}^*} \left(\frac{2}{p}\right) \right)^{-\frac{1}{2}}. \]

We now fix the parameters as functions of $x$ as follows:
\[ r(x) := \lfloor \log \log \log x \rfloor, \quad y(x) := \frac{\log x}{3 \log \log x}, \quad \epsilon(x) := \frac{1}{13y(x)}. \]

The explicit asymptotic evaluation of each term as $x \to \infty$ is as follows:

For the Petrov term, given the choice of $\epsilon$, the upper bound for the primes in the effective set is exactly $y$. Thus, $\mathcal{P}^* = \{5 \leq p \leq y\}$. Since excluding a few constant primes does not affect the asymptotic behavior, Mertens' theorem yields:
\[ \sum_{p \in \mathcal{P}^*} \left(\frac{2}{p}\right) \sim 2 \log \log y \sim 2 \log \log \log x \]
Thus, this term is of order $O((\log \log \log x)^{-1/2})$.

For the $E_1$ term, substitution and simplification yield $\frac{13}{\log y} \sim \frac{13}{\log \log x}$. On the other hand, $(\log \log x)^{-1} = o((\log \log \log x)^{-1/2})$. Hence, this term is asymptotically smaller than the Petrov term.

For the $E_2$ term, we have $2^{-r} \asymp (\log \log x)^{-\log 2}$. Since  $\log 2 > 0.5$ and $(\log \log x)^{-\log 2} = o((\log \log \log x)^{-1/2})$, this term is also dominated by the Petrov term.

For the periodic error term, since $ry = o(\log x)$, for sufficiently large $x$ we always have $\exp(1.02ry) \leq x^{1/2}$. Consequently, this term is bounded by $x^{-1/2}$ and thus by $o((\log \log \log x)^{-1/2})$.

For the $E_0$ term, since it equals $\frac{13y}{x}$, it is strictly $o(1)$ compared to the logarithmic terms.

Since all error terms are strictly of an order smaller than $(\log \log \log x)^{-1/2}$, they are all absorbed by the Petrov term. Consequently, the asymptotic upper bound is exclusively governed by the behavior of the concentration function, and we strictly obtain:
\[ A_k(x) \ll_k \frac{x}{\sqrt{\log \log \log x}} \]

This concludes the proof.
\end{proof}

\end{document}
